%--------------------------------------------------------------------------------
%
%
%--------------------------------------------------------------------------------
\pagebreak
\section*{NM - create a new T64 module}
\addcontentsline{toc}{section}{NM - create a new T64 module}

%--------------------------------------------------------------------------------
\begin{iPageDescEntry}{Syntax}
\texttt{NM <mType> , <modNum>  [ \{ , <key>=<val> \}* ] } \\
\end{iPageDescEntry}

%--------------------------------------------------------------------------------
\begin{iPageDescEntry}{Description}

A Twin-64 system is a set of modules with at least one processor, one memory and one I7O module. The \texttt{NM} command adds a module to the Twin-64 simulator module table. The command is intended to be used by a configuration file. A configuration file will simply contains a list of commands to setup a system. 

The \texttt{module type} parameter defines the module to be added. There are three types, \texttt{PROC}, \texttt{MEM} and \texttt{IO}. The module number parameter assigns the module number. The Twin-64 system supports up to 64 modules at the system level. 

The remaining optional parameters pass configuration values for the module. They all are of the form "\texttt{name=val}". 

\begin{flushleft}
\begin{tabularx}{\linewidth}{|l|l|X|}
\hline
\textbf{Keyword}  & \textbf{Value} & \textbf{Usage} \\ \hline
\texttt{SPA\_ADR} & \texttt{32-bit address} & Soft physical address start. The address must be aligned to the length parameter. \\ \hline
\texttt{SPA\_LEN} & \texttt{32-bit size} & Soft physical address length. The size is the pageSize multiplied by a power of 2.\\ \hline
\texttt{TLB} & \texttt{TLB\_FA\_xxS} & \ TLB set size. The defined values are 16, 32, 64 and 128. \\ \hline
\texttt{MEM} & \texttt{RAM | ROM} & Memory access type. There are read/write and read only memories. \\ \hline
\end{tabularx}
\end{flushleft}

\end{iPageDescEntry}

%--------------------------------------------------------------------------------
\begin{iPageDescOpEntry}{Example}
\begin{lstlisting}[style=iPageOpStyle]

# add a processor module - 3 - with defaults.
NM proc, 3

# add a processor module - 4 - with a larger TLB                   
NM proc, 4, tlb=TLB_FA_128S

# add a memory module covering the first 256 pages
NM mem, 7, spa_adr=0x0, spa_len=256*0x1000

\end{lstlisting}
\end{iPageDescOpEntry}

%--------------------------------------------------------------------------------
\begin{iPageDescEntry}{Notes}
None.
\end{iPageDescEntry}
